\documentclass[11pt,letterpaper]{article}
\PassOptionsToPackage{hyphens}{url}
\usepackage[margin=.83in,top=.8in,bottom=.8in,headheight=15pt]{geometry}
\usepackage[T1]{fontenc}
\usepackage{lmodern,amsmath,amssymb,mathtools,bm,booktabs,array,graphicx,tikz}
\usetikzlibrary{arrows.meta,positioning,calc,decorations.pathreplacing,patterns}
\definecolor{ink}{HTML}{17344D}\definecolor{accent}{HTML}{087E8B}\definecolor{pale}{HTML}{EDF5F5}\definecolor{gold}{HTML}{AD6B19}
\providecommand{\dd}{\,\mathrm{d}}\providecommand{\R}{\mathbb{R}}
\providecommand{\E}{\mathbb{E}}\providecommand{\Var}{\operatorname{Var}}
\providecommand{\Cov}{\operatorname{Cov}}\providecommand{\diag}{\operatorname{diag}}
\providecommand{\norm}[1]{\left\lVert#1\right\rVert}
\providecommand{\abs}[1]{\left|#1\right|}
\usepackage{microtype,hyperref,fancyhdr,enumitem,amsthm,sectsty,needspace}
\allsectionsfont{\raggedright}
\hypersetup{colorlinks=true,linkcolor=ink,urlcolor=accent,pdfauthor={Zan Ahmad},pdftitle={Linear observations and a two-compartment exchange model / Review}}
\setlength{\parindent}{0pt}\setlength{\parskip}{6pt}\linespread{1.07}
\setlist{itemsep=3pt,topsep=4pt,leftmargin=*}
\emergencystretch=2em\clubpenalty=10000\widowpenalty=10000
\pagestyle{fancy}\fancyhf{}\fancyhead[L]{\small Scientific Computing in Python / 2027}\fancyhead[R]{\small Lecture 3}
\fancyfoot[L]{\footnotesize Zan Ahmad}\fancyfoot[R]{\footnotesize\thepage}\renewcommand{\headrulewidth}{.3pt}
\newtheorem{proposition}{Proposition}\newtheorem{theorem}{Theorem}\newtheorem{lemma}{Lemma}
\theoremstyle{definition}\newtheorem{definition}{Definition}\newtheorem{example}{Example}
\numberwithin{equation}{section}
\begin{document}
\setlength{\parskip}{4pt}\setlist{itemsep=2pt,topsep=3pt,leftmargin=*}
\section*{Linear observations and a two-compartment exchange model --- review sheet}
Use this after attempting the worked examples and practice. Explain each claim in words before calculating.

\section*{Key ideas}
\begin{itemize}
\item The observation map and state evolution are different equations.
\item One transfer creates equal-and-opposite mass rates.
\item The sum is conserved while the difference decays.
\item Conservation does not measure individual-state integration error.
\item Fitting requires data informative about the parameter.
\item Verification, fit and scientific validation support different claims.
\end{itemize}
\section*{Common errors and repairs}
\textbf{Use k rather than 2k in the difference decay.} Both masses change toward each other, so subtracting derivatives gives -2k times the difference.

\textbf{Call t\_eval the fixed integration step.} It specifies returned times; adaptive internal steps are separate.

\textbf{Treat low fit residual as validation of the mechanism.} It measures agreement under the selected model and objective only.

\section*{Retrieve without your notes}
\begin{itemize}
\item What is the equilibrium for a closed equal-volume pair totaling 14 mg?
\item What happens at k=\allowbreak{}0?
\item What units does k h have?
\end{itemize}
\section*{Readiness check}
\begin{itemize}
\item Recover two initial masses from a stated linear sensor map and check its residual.
\item Translate one transfer rate into equal-and-opposite state derivatives with units.
\item Compare an ODE solver and Euler steps against an exact solution and mass conservation.
\item Fit one exchange parameter to synthetic observations while distinguishing fit residual from integration error.
\item Identify a data configuration that cannot determine the exchange rate.
\end{itemize}
\end{document}
